\documentclass[11pt]{article}
\usepackage{mathblatt}
% Präambel des Setzers (werkzeuge/setzer-praeambel.tex, wortgleich); dazu \gteil (Muster T2U)
\makeatletter
\setlength{\mbpfbe}{0mm}% keine Punkte (pruefung.md § 5: nur beim Prüfstein)
\newlength{\gEin}\setlength{\gEin}{0pt}
\let\gSchrift\relax
\renewcommand{\pfkreuzzeile}[1]{\par\noindent\hangindent3.2em\hangafter1\hspace*{1.6em}$\square$\ #1\par}
\newcommand{\gkreuz}[1]{$\square$\ #1\hspace{2em}}
\newcommand{\gkreuzzeile}[2]{\par\noindent\hangindent2.2em\hangafter1\mbox{#1\,$\square$}\ #2\par}% Buchstabe, Option: eine je Zeile (neben Skizze)
% links, was man ansieht, rechts, was man tut; Spaltengrenze überall an derselben Stelle (gemeinsam.md „Satz“)
\newsavebox{\gbox}\newlength{\gLinks}
% v0.10 (Lehrer 13.10.): keine starre Halbierung – die linke Spalte ist so breit wie die Figur (+ 6 mm), die
% Antwort steht direkt daneben; ist die Figur breiter als 0,62 Zeile, steht die Antwort darunter
\newcommand{\gzwei}[2]{\par\smallskip\noindent\sbox{\gbox}{#1}%
  \setlength{\gLinks}{\dimexpr\wd\gbox+6mm\relax}%
  \ifdim\gLinks>0.62\linewidth
    \usebox{\gbox}\par\smallskip\noindent\begin{minipage}{\linewidth}#2\end{minipage}%
  \else
    \begin{minipage}[t]{\gLinks}\vspace{0pt}\usebox{\gbox}\end{minipage}%
    \begin{minipage}[t]{\dimexpr\linewidth-\gLinks\relax}\vspace{0pt}\raggedright #2\end{minipage}%
  \fi\par}
% v0.9 (12.10., pruefung.md § 5 „Marke mit B, Bedingungen rechts“): drittes Argument von pfaufg = Bedingungen
% der Prüfung (v0.10: \pfohnetr = TR durchgestrichen, „FOR“), klein grau im rechten Rand auf der Grundlinie der Nummer, Zeilen mit \\
\newcommand{\gbedrand}[2]{\ifblank{#2}{}{\rlap{\hspace*{\dimexpr#1+2mm\relax}%
  \parbox[t]{15mm}{\raggedright\scriptsize\leavevmode\color{mbgrau}\strut#2}}}}
% Teil eines Bündels (a, b …): Marke links im Rand, Bedingungen rechts im Rand, auf der Zeile des Teils
\newcommand{\gteilrand}[2]{\llap{\mbpfrandmarke{#1}\hspace*{\mbpfnr}}\gbedrand{\linewidth}{#2}}
% Teil eines Bündels: Buchstabe in eigener Spalte, Inhalt eingerückt (\gEin = Einzug, damit die Spaltengrenze von
% \gzwei an derselben Stelle bleibt)
\newenvironment{gbteil}[3]{\par\smallskip\noindent\gteilrand{#2}{#3}\makebox[\mbpfteil][l]{\textbf{#1}}%
  \begin{minipage}[t]{\dimexpr\linewidth-\mbpfteil\relax}\setlength{\gEin}{\mbpfteil}\raggedright\strut\ignorespaces}%
 {\end{minipage}\par}
% Grafik höchstens so breit wie die Spalte (Bündel nebeneinander, Koordinatensystem ohne Karo)
\newsavebox{\gpbox}
\newcommand{\gpassend}[1]{\sbox{\gpbox}{#1}\ifdim\wd\gpbox>\linewidth\typeout{GPASSEND: verkleinert \the\wd\gpbox\space auf \the\linewidth}\resizebox{\linewidth}{!}{\usebox{\gpbox}}\else\usebox{\gpbox}\fi}
% Skizzenfeld (v0.10, Lehrer 13.10.): Karo ab dem linken Rand über die Zeile, „Skizze“ klein grau oben links im Feld
% (wie \pfskizzenfeld der Vorlage); der rechte Rand bleibt den Bedingungen
\newcommand{\gskizze}[1]{\par\smallskip\noindent\begin{tikzpicture}
  \clip (0,0) rectangle ({\linewidth-1mm},{#1*10mm});
  \draw[black!16,line width=0.3pt,step=5mm] (0,0) grid ({\linewidth},{#1*10mm});
  \node[anchor=north west,inner sep=2pt,font=\scriptsize,text=mbgrau,fill=white] at (0,{#1*10mm}) {Skizze};
  \end{tikzpicture}\par}
\RenewDocumentEnvironment{pfaufg}{m m m}{%
  \begin{lrbox}{\mbaufgabebox}\begin{minipage}[t]{\linewidth}%
  \gSchrift\noindent\hspace*{\gEin}\mbpfrandmarke{#1}\gbedrand{\linewidth-\gEin}{#3}%
  \begin{minipage}[t]{\mbpfnr}\strut\textbf{#2}\end{minipage}%
  \begin{minipage}[t]{\dimexpr\linewidth-\gEin-\mbpfnr\relax}\raggedright\strut\ignorespaces}%
 {\end{minipage}%
  \end{minipage}\end{lrbox}%
  \setlength{\mbaufgabehoehe}{\dimexpr\ht\mbaufgabebox+\dp\mbaufgabebox\relax}%
  \ifdefempty{\gkopfnach}{\par\addvspace{5pt}\Needspace*{\mbaufgabehoehe}}{%
    % Stufenkopf, der mit dieser (ersten) Aufgabe zusammenbleibt: Platz für beide verlangen, dann den Kopf setzen
    % (12.10., Quadratische: ein Kopf stand allein unten, die Aufgabe mit Koordinatensystem auf der nächsten Seite)
    \par\Needspace*{\dimexpr\mbaufgabehoehe+40pt\relax}\gkopfnach\gdef\gkopfnach{}\par\addvspace{5pt}}%
  \noindent\usebox{\mbaufgabebox}\par\addvspace{5pt}}
% Stufenkopf vormerken: \gkopfvor{<Kopf>} – die nächste pfaufg setzt ihn, sobald beide zusammen Platz haben
\def\gkopfnach{}
\newcommand{\gkopfvor}[1]{\gdef\gkopfnach{#1}}
% Teilaufgabe mit Skizze links in gzwei: Buchstabe in eigener Spalte, darunter Text und Skizze
\newcommand{\gteil}[2]{\makebox[\mbpfteil][l]{\textbf{#1}}\parbox[t]{\dimexpr0.5\textwidth-\gEin-\mbpfnr-5mm-\mbpfteil\relax}{\raggedright #2}}
% Stufenkopf der Vorlage (mathblatt.sty 2026-10-12: Linie, Abstand, 13 pt, \pfjahre) ohne ihren Vorhalt von
% 14 Zeilen: der Setzer fordert vor dem Kopf 0,16 Seitenhöhe (Kopf + erste Aufgabe), mehr nicht
\setlength{\mbpfstufevorhalt}{0pt}
\makeatother
\newcommand{\karo}[1]{\par\smallskip\noindent\begin{tikzpicture}
  \clip (0,0) rectangle ({\linewidth-1mm},{#1*10mm});
  \draw[black!16,line width=0.3pt,step=5mm] (0,0) grid ({\linewidth},{#1*10mm});
  \end{tikzpicture}\par}
\newcommand{\kk}[1]{$\square$\,#1\hspace{0.9em}}
\newcommand{\linie}[1]{\,{\color{black!45}\rule[-1pt]{#1}{0.4pt}}\,}
\newcommand{\pkt}{\textperiodcentered{}}
% v0.14 (14.10.): Stufenkopf mit fünf Punkten (Option = Jahre für die Bestellkarte) oder „zuletzt 19“
\NewDocumentCommand{\gpunkte}{o m}{#2}
\newcommand{\gzuletzt}[1]{zuletzt #1}
% v0.7 (12.10.): Schülerfrage unter dem Stufenkopf; Ankreuzoption als Bild (Buchstabe und Kästchen darunter)
\newcommand{\gfrage}[1]{\par\nobreak\noindent\hspace*{\mbpfnr}{\small\itshape #1}\par\nobreak\addvspace{2pt}}
\newlength{\gbildbreite}\setlength{\gbildbreite}{0.32\linewidth}% Setzer setzt sie je Reihe (vier Bilder: 0,24)
\newcommand{\gbildoption}[3]{\begin{minipage}[t][#3][b]{\gbildbreite}\centering #2\par\smallskip\mbox{#1\,$\square$}\end{minipage}}
% Lineare (12.10.): Bild mit Buchstabe ohne Kästchen (Diagramm zuordnen); Aussage mit „wahr / falsch“ zum Ankreuzen
\newcommand{\gbildname}[3]{\begin{minipage}[t][#3][b]{\gbildbreite}\centering #2\par\smallskip\mbox{\textbf{#1}}\end{minipage}}
\newcommand{\gwahrfalsch}[2]{\par\noindent\hangindent2.2em\hangafter1\makebox[2.2em][l]{(#1)}#2\hfill\mbox{wahr $\square$\quad falsch $\square$}\par}
\begin{document}

\pfheftstil
\pflaufkopf[Abitur GK \pkt{} Prüfung]{Integral berechnen \pkt{} Lösungen}{AA140}\pfseitevon
\pfheftkopf{Integral berechnen}{Lösungen}
\pfstufenzeile{Stammfunktion nachweisen \pkt{} Integral gesucht \pkt{} Grenze oder Parameter gesucht}

\Needspace*{0.16\textheight}\pfstufekopf{Stammfunktion nachweisen}{\gpunkte[24 23 22]{\textcolor{black!25}{$\circ$}\,\textcolor{black!25}{$\circ$}\,\textcolor{black!80}{$\bullet$}\,\textcolor{black!80}{$\bullet$}\,\textcolor{black!80}{$\bullet$}}}\special{pdf:dest (stufe-stammfunktion-nachweisen) [@thispage /XYZ @xpos @ypos null]}
\begin{pfloesungb}{}
\lzb{1.}{}{F' = f(x)}{F'(x) = $-$3 · $e^{ - x/3 + 2}$ + ($-$3x) · ($-\tfrac{1}{3}$) · $e^{ - x/3 + 2}$ = (x $-$ 3) · $e^{ - x/3 + 2}$ = f(x)}
\end{pfloesungb}
\begin{pfloesungb}{}
\lzb{2.}{}{F(x) = ($\tfrac{2}{3}$ x² $-$ $\tfrac{8}{3}$ x $-$ $\tfrac{2}{3}$) · $e^{x/2}$; A $\approx$ 1457 cm²}{F(x) = ($\tfrac{2}{3}$ x² $-$ $\tfrac{8}{3}$ x + $\tfrac{16}{3}$) · $e^{x/2}$ $-$ 6$e^{x/2}$ = ($\tfrac{2}{3}$ x² $-$ $\tfrac{8}{3}$ x $-$ $\tfrac{2}{3}$) · $e^{x/2}$; Integral von $-$3 bis 0 = F(0) $-$ F($-$3) = $-\tfrac{2}{3}$ $-$ $\tfrac{40}{3}$ · $e^{ - 1{,}5}$ $\approx$ $-$3,642 $\Rightarrow$ A = 4 · 3,642 $\approx$ 14,57 LE²; 1 LE² = 100 cm² $\Rightarrow$ A $\approx$ 1457 cm²}
\end{pfloesungb}
\begin{pfloesungb}{}
\lzb{3.}{}{F' = f(x); gleicher Inhalt}{F'(x) = $-\tfrac{1}{18}$ x³ + $\tfrac{1}{3}$ x² = f(x); F(8) $-$ F(0) = 0 $\Rightarrow$ Integral von 0 bis 8 ist 0; f $\ge$ 0 auf [0; 6], f $\le$ 0 auf [6; 8] $\Rightarrow$ die beiden Flächenstücke heben sich auf $\Rightarrow$ gleicher Inhalt}
\end{pfloesungb}

\Needspace*{0.16\textheight}\pfstufekopf{Integral gesucht}{\gpunkte[26 25 23]{\textcolor{black!80}{$\bullet$}\,\textcolor{black!80}{$\bullet$}\,\textcolor{black!25}{$\circ$}\,\textcolor{black!80}{$\bullet$}\,\textcolor{black!25}{$\circ$}}}\special{pdf:dest (stufe-integral-gesucht) [@thispage /XYZ @xpos @ypos null]}
\begin{pfloesungb}{}
\lzb{4.}{}{F(0) $-$ F($-$1) = $\tfrac{23}{12}$}{F(x) = $\tfrac{1}{4}$ x$^{4}$ + $\tfrac{2}{3}$ x³ $-$ $\tfrac{3}{2}$ x² $\Rightarrow$ F(0) $-$ F($-$1) = 0 $-$ ($\tfrac{1}{4}$ $-$ $\tfrac{2}{3}$ $-$ $\tfrac{3}{2}$) = $\tfrac{23}{12}$}
\end{pfloesungb}
\begin{pfloesungb}{}
\lzb{5.}{}{F(2) $-$ F($-$2) $\approx$ 13,97; Abweichung $\approx$ 5,8 \%}{F(2) $-$ F($-$2) = 2e² $-$ 6$e^{ - 2}$ $\approx$ 13,97; Näherung 2e² $\approx$ 14,78 $\Rightarrow$ (14,78 $-$ 13,97)/13,97 $\approx$ 0,058 $\Rightarrow$ Abweichung $\approx$ 5,8 \%}
\end{pfloesungb}
\begin{pfloesungb}{}
\lzb{6.}{}{$\tfrac{875}{5}$ = 175; in den ersten fünf Minuten beträgt die Temperatur im Mittel 175 °C}{Integral von 0 bis 5 = [x$^{4}$ $-$ 15x³ + 81x² + 20x] von 0 bis 5 = 875 $\Rightarrow$ $\tfrac{875}{5}$ = 175; in den ersten fünf Minuten beträgt die Temperatur im Mittel 175 °C}
\end{pfloesungb}

\Needspace*{0.16\textheight}\pfstufekopf{Grenze oder Parameter gesucht}{\gpunkte[26 24 22]{\textcolor{black!80}{$\bullet$}\,\textcolor{black!25}{$\circ$}\,\textcolor{black!80}{$\bullet$}\,\textcolor{black!25}{$\circ$}\,\textcolor{black!80}{$\bullet$}}}\special{pdf:dest (stufe-grenze-oder-parameter-gesucht) [@thispage /XYZ @xpos @ypos null]}
\begin{pfloesungb}{}
\lzb{7.}{}{a $\approx$ 0,31}{Integral von $-$1 bis 1 über f = [0,1x$^{5}$ $-$ x³ + 2,5x] von $-$1 bis 1 = $\tfrac{16}{5}$ $\Rightarrow$ a · $\tfrac{16}{5}$ = 1 $\Rightarrow$ a = $\tfrac{5}{16}$ $\approx$ 0,31}
\end{pfloesungb}
\begin{pfloesungb}{}
\lzb{8.}{}{Schnittstellen m und 0; m = $-$3}{Schnittstellen m und 0; Integral von m bis 0 über (mx $-$ x²) = $-$m³/6 = $\tfrac{9}{2}$ $\Rightarrow$ m³ = $-$27 $\Rightarrow$ m = $-$3}
\end{pfloesungb}
\begin{pfloesungb}{}
\lzb{9.}{}{= $\tfrac{3}{2}$; a $\approx$ $-$0,47}{Gesamtfläche 2 $-$ 2$e^{ - ln 4}$ = 2 $-$ $\tfrac{1}{2}$ = $\tfrac{3}{2}$; Fläche von a bis 0: 2 $-$ 2$e^{a}$ = $\tfrac{3}{4}$ $\Rightarrow$ $e^{a}$ = $\tfrac{5}{8}$ $\Rightarrow$ a = ln($\tfrac{5}{8}$) $\approx$ $-$0,47}
\end{pfloesungb}
\end{document}